\chapter{}

Para tentar equilibrar nossa situação, foi preciso vender um apartamento que me coubera numa primeira partilha da herança do meu pai.
Não vou negar, a herança significava para mim a certeza de que eu teria como sustentar a faculdade dos meninos.
Não podia dizer isso ao Paulo sem ofendê-lo mortalmente, mas era verdade.
Entretanto, era vender o apartamento, ou pedir emprestado a parentes, ou aos bancos.
Tanto minha família, como meu sogro, tinham nos emprestado dinheiro, mas eu não queria que aquilo continuasse, até para proteger Paulo.
Conhecia a cabeça deles todos e não ia aguentar qualquer insinuação de que estivéssemos nos aproveitando da situação para parasitar quem quer que fosse.

Hoje, sem nunca ter conseguido abrir-me a respeito do assunto com Paulo, acabei chegando à conclusão de que tamanha sempre foi a insegurança dele quanto à própria capacidade de garantir nossa subsistência, que suas reações nunca eram proporcionais à realidade, mas ao medo que ele sentia de não dar conta do recado.
Ele sempre achava, como ainda acha, que precisava ganhar muito dinheiro e negócios para ganhar muito dinheiro, é claro, exigiam um investimento que nunca tínhamos condições de bancar, era o que eu pensava.
Porém, Paulo achava que os bancos estavam aí era para isso mesmo.

A essa forma de pensar, juntava-se uma outra característica que eu também nunca tinha podido aceitar: Paulo achava que o lance para ganhar dinheiro era pensar em coisas incomuns, novas, inéditas, para o que concorria, é claro, sua eterna busca por reconhecimento.
Eu até concordava, em parte, desde que tais coisas fossem ditadas por uma necessidade bem concreta, constatada e ainda não satisfeita.
Mas essa discordância era outro pé para discussões intermináveis que acabavam inevitavelmente em briga.
E tudo que eu não queria era recriar em casa o clima que nós dois havíamos enfrentado em nossas próprias famílias.
Eu preferia calar, mas aquilo estava me fazendo muito mal.
Eu me sentia covarde e impotente por não enfrentá-lo.

Depois de muito procurar aqui e ali alguma coisa para ganhar dinheiro, Paulo se entusiasmou com a ideia de criar camarões da Malásia.
Já se comentava que os camarões marinhos viriam a faltar dentro de alguns anos dada a pesca predatória e as alterações observáveis nos oceanos.
Alguns biólogos começavam a estudar esse gigante da Malásia como uma alternativa interessante para ser criada em cativeiro em nosso país.
E ao final do seu trabalho na SULANOR, Paulo já tinha montado no seu sítio de Miguel Alves um cultivo experimental e com bastante sucesso.
Mas, os alevinos tinham que vir de longe, com dificuldade e a um alto preço.
E, não raro, chegavam mortos.
Então ele decidiu que ia produzir os alevinos num laboratório a ser construído inicialmente em nosso quintal.
E, a ideia era fazê-lo inspirando-se nas instalações simples e primitivas usadas pelos produtores do extremo oriente.


\begin{figure}
\centering
\includegraphics[width=0.89\linewidth]{38/+paulo-na-criação.png}
\caption{Com o Paulo, na criação de camarões da Malásia em Miguel Alves-PI.}
\end{figure}

\begin{figure}
\centering
\includegraphics[width=0.89\linewidth]{38/eu+gigante.png}
\caption{Eu e o gigante da Malásia.}
\end{figure}

Quando Paulo mete uma coisa na cabeça, torna-se compulsivo.
Até que esse laboratório ficasse pronto, ele trabalhou até o limite da exaustão.
Quando finalmente começou a funcionar, meses depois, eu não acreditava em meus próprios olhos.
Era, embora rudimentar quanto aos materiais, uma coisa extremamente complexa, uma maçaroca de fios e canos, alguns capilares quase, descendo do teto sobre os tanques, levando água e oxigênio a um monte de compartimentos de tamanhos e funções variadas, todos borbulhando à espera dos minúsculos ovinhos.

A notícia da existência do laboratório provocou grande repercussão entre os pesquisadores locais e até de fora.
Um dia, parou diante do nosso portão uma excursão de sorridentes japoneses em trajes de turista, câmeras fotográficas a tiracolo que, entre profusas mesuras, fizeram saber, através do intérprete, que queriam fotografar o laboratório para mostrá-lo em seu país, do outro lado do mundo.
Com o tempo, nos habituaríamos a ver algumas celebridades aparecerem em nosso quintal.

Outra façanha do Paulo foi treinar o pessoal que trabalharia com ele.
Pude avaliar a extensão da dificuldade quando a mais preparada das estagiárias, a que nos fora indicada como a melhor aluna da série final do curso de Biologia da Universidade Federal do Piauí, descreveu-me sua observação de um camarão ao microscópio do qual, minutos antes, eu levara as lentes para lavar e ainda não trouxera de volta.
Apesar disso, Paulo conseguiu tocar o laboratório e produzir os filhotes, tendo como seu braço direito uma moça que ele tirara da enxada e duas ou três das estagiárias.

O dinheiro acabara de novo e era preciso pensar nos tanques de engorda.
De novo iríamos aos bancos.
O gerente do Banco do Nordeste ofereceu ao Paulo repassar-lhe uma fazenda tomada a um proprietário inadimplente que, de lambuja, trazia acoplada uma área de mil e quatrocentos hectares no sul do Piauí que o sujeito dera em garantia ao banco e que poderia servir de novo a essa mesma finalidade.
Paulo achou um bom negócio, assumiu a fazenda e tomou então um empréstimo para a construção dos tanques.

Pôr de novo a Coordenação a funcionar deu-me um bocado de trabalho.
De modo que acompanhei de longe a fase da implantação do projeto “Camarão Azul”, a fazenda de aquicultura do Paulo que começava a tomar forma.
Vinte e nove grandes tanques abastecidos por água bombeada do rio Parnaíba foram construídos.
Provisoriamente, o laboratório fora transferido de casa para uma instalação abandonada do Estado.
Nosso espaço do quintal era insuficiente para produzir os filhotes necessários.
Porém, o que Paulo tinha em mente era transferi-lo de vez para a fazenda.
Para isso, um grande galpão começou a ser construído logo à entrada da propriedade.

Nem por um momento duvidei que Paulo fosse ter sucesso na implantação do projeto.
Nossos conhecidos comentavam espantados que a fazenda do Paulo podia ser vista pela janela dos aviões de carreira.
Em termos de Piauí, o porte do empreendimento causava admiração.
Contudo, conhecendo o que ele fizera antes, eu sabia que para ele aquilo era quase nada.
E esse era o maior problema.
A capacidade de realização do Paulo ultrapassava em muito os recursos que sozinhos conseguíamos arrecadar.
Estávamos encalacrados.
Foi então que surgiu um empresário muito rico, interessado no negócio, e que se propôs como sócio capitalista.

Esse senhor, de cujo nome não me lembro, enriquecera como “coureiro”, ou caçador de jacarés do Pantanal e depois se estabelecera com plantio e beneficiamento de arroz no Rio Grande do Sul.
Era um sujeito com uma aparência rude de quem trabalhou duro para conquistar seu quinhão.
Na verdade, havíamos de descobrir com o tempo, que o que ele e o irmão que o acompanhava queriam era arrumar uma ocupação séria para os filhos, dois rapazinhos bonitinhos e obviamente familiarizados com todos os modismos em voga, mas pouco afeitos ao trabalho.
Mas, o importante é que aceitando os dois jovens como sócios, Paulo conseguiria pôr o projeto a funcionar.
E assim foi feito.

Paulo produziu camarões da Malásia, talvez os maiores e mais bonitos que o mercado viu, a julgar pela reação dos clientes do Rio e de São Paulo para quem enviávamos semanalmente, por avião, grandes caixas de isopor cheias.
Mas o produto exigia cuidados extremos no transporte e no preparo, sem o que adquiria uma consistência desagradável e um sabor amargo.
Não conseguiu suplantar os camarões tradicionais na preferência das donas de casa.
E seu maior trunfo, o tamanho, não podia ser explorado porque sua manutenção nos tanques era cara, o que determinava que ele fosse despescado ainda pequeno.
Eu, que o cozinhava em casa, aos montes (embora apertados, comemos camarão como nunca), percebi que, comercialmente, o gigante da Malásia estava fadado ao fracasso.
E realmente, em pouco tempo, Paulo e outros produtores tiveram que começar a procurar alternativas para não perder o investimento feito.
A Camarão Azul passaria a produzir peixes.

Em conversa com os jovens sócios, surgiu a ideia de vender peixes-vivos em tanques espalhados pelos mercados e supermercados da cidade.
A iniciativa teve enorme sucesso, o que aparentemente despertou a ambição dos dois rapazes, especialmente do filho do coureiro.
Eles chegaram ao Paulo com a proposta de desmanchar a sociedade.
Provavelmente, achando os peixes bem menos complicados que os camarões decidiram que, daí por diante, seguiriam sozinhos.

Ocorre que durante todo esse processo, quando foi preciso novamente recorrer aos bancos, Paulo, para não perder participação no negócio, ofereceu em garantia a única coisa que nos tinha sobrado: a nossa casa.
Aquilo me arrasou.
E o pior é que nesse momento em que a sociedade entrava em dissolução, o filho do coureiro parecia determinado a ficar com tudo, inclusive a casa, alegando que o Paulo, até então, só entrara com o trabalho.
Chegou a me telefonar fazendo ameaças veladas, caso nos recusássemos a entregá-la.
Eu, para azar dele, andava de tal forma esgotada, irritada, cansada de tudo, que perdi as estribeiras e disse-lhe ao telefone tudo quanto os pais dele lhe deviam ter dito quando ainda era tempo.
Ele emudeceu.
Felizmente a lei nos salvou.
Às vezes isso acontece no Brasil.
O Banco não podia ter aceitado a nossa casa em garantia.
Era a única que tínhamos.

Nessa altura dos acontecimentos, Tomaz, Marcelo e Fernando já estavam em São Paulo, os três cursando universidades públicas.
Ainda assim, tínhamos dificuldades em mandar-lhes dinheiro todo mês.
Tomaz conseguira bolsa, Marcelo dava aulas em cursinho, Fernando estudava em período integral e agora chegava a hora de Paula também prestar vestibular.
Ao mesmo tempo, em Araraquara, com a maioridade do caçula, Paulinho, a partilha da herança deixada pelo comandante voltou à baila.

Se as coisas para nós não corriam bem, o restante da família, tanto do lado do Paulo como do meu, vivia tragédias com certeza piores.
Paulo perdeu seu irmão Inácio que tinha com ele uma ligação muito forte e João, meu irmão, perdeu a mulher que idolatrava, num parto mal conduzido, do qual lhe ficou nos braços um filho recém-nascido e uma revolta duríssima de aplacar.

A morte da Bel cindiu a vida de João ao meio e marcou-o para sempre.
Por meses, ele errou sem rumo, sem achar lugar, agarrado ao filho recém-nascido que ele não entregou a ninguém, descrente de Deus, da vida, da sua profissão de médico e preso de uma raiva que ele arremessava sobre tudo e todos, sem muito critério.
E, naquele momento o alvo mais vistoso e disponível era Reginaldo, nosso irmão mais velho, que protagonizaria, no papel de vilão, o drama em que se transformou a partilha do restante dos bens deixados por nosso pai.

Algum tempo depois, nova tragédia se abateria sobre nossa família: Renatão, após longo sofrimento, morreu, aos cinquenta e dois anos, vítima de câncer no fígado.

\begin{figure}[H]
\centering
\includegraphics[width=0.9\linewidth]{38/renato+fernando.png}
\caption{Renatão e Fernando, inseparáveis, mestre e discípulo na devoção ao churrasco.}
\end{figure}